\chapter{Oncology}

\medterm{Neoplasm}-- Abnormal growth of tissue, may be benign or malignant. Benign neoplasms: localized, slow-growing, do not metastasize. Malignant neoplasms (cancer): invasive, metastasize, can be fatal. Classification: carcinoma (epithelial origin), sarcoma (mesenchymal origin), leukemia (hematopoietic), lymphoma (lymphoid).

\medterm{Carcinoma}-- Malignant tumor of epithelial origin. Types: squamous cell carcinoma, adenocarcinoma, basal cell carcinoma, transitional cell carcinoma. Most common cancer type (~80-90% of malignancies). Examples: lung carcinoma, breast carcinoma, colon carcinoma.

\medterm{Sarcoma}-- Malignant tumor of mesenchymal origin (bone, cartilage, fat, muscle, blood vessels). Types: osteosarcoma, chondrosarcoma, liposarcoma, rhabdomyosarcoma, angiosarcoma. Less common than carcinomas (~1% of adult cancers).

\medterm{Leukemia}-- Malignant proliferation of hematopoietic cells in bone marrow. Types: acute lymphoblastic (ALL), acute myeloid (AML), chronic lymphocytic (CLL), chronic myeloid (CML). Acute: rapid onset, immature cells (blasts). Chronic: slower progression, mature cells.

\medterm{Lymphoma}-- Malignant tumor of lymphocytes. Types: Hodgkin lymphoma (Reed-Sternberg cells), non-Hodgkin lymphoma (diverse subtypes). Hodgkin: more common in young and older adults, predictable spread. NHL: more common overall, variable behavior.

\medterm{Tumor}-- Abnormal mass of tissue, may be benign or malignant. Tumors classified by tissue of origin. Terminology: -oma usually benign, but exceptions (melanoma, lymphoma, seminoma are malignant).

\medterm{Metastasis}-- Spread of cancer cells from primary site to distant locations. Routes: lymphatic, hematogenous, direct extension, transcoelomic. Common sites: lungs, liver, bone, brain. Distant metastasis indicates Stage IV disease.

\medterm{Malignancy}-- See cancer.

\medterm{Benign Tumor}-- Non-cancerous growth, slow-growing, well-differentiated, localized, does not invade or metastasize. Examples: lipoma, fibroma, adenoma. May cause problems by compression or hormone secretion.

\medterm{Malignant Tumor}-- Cancerous growth, invasive, can metastasize, poorly differentiated. Requires aggressive treatment.

\medterm{Cancer Staging}-- Classification of cancer extent. TNM system: T (tumor size/extent), N (nodal involvement), M (metastasis). Stages: 0 (in situ), I (early), II-III (advanced), IV (metastatic). Guides treatment and prognosis.

\medterm{Grade}-- Degree of cellular differentiation. Grade 1: well-differentiated. Grade 2: moderately differentiated. Grade 3: poorly differentiated. Grade 4: undifferentiated. Higher grade = more aggressive.

\medterm{Biopsy}-- Tissue sampling for diagnosis. Types: needle biopsy, excisional biopsy, incisional biopsy, endoscopic biopsy. Essential for cancer diagnosis and treatment planning.

\medterm{Chemotherapy}-- Use of cytotoxic drugs to treat cancer. Mechanisms: DNA damage, microtubule inhibition, antimetabolite action. Common agents: alkylating agents, anthracyclines, taxanes, platins, antimetabolites. Side effects: myelosuppression, nausea, alopecia.

\medterm{Radiation Therapy}-- Use of ionizing radiation to kill cancer cells. Types: external beam, brachytherapy, stereotactic. Mechanisms: DNA damage, cell death. Side effects: fatigue, skin reactions, site-specific effects.

\medterm{Immunotherapy}-- Treatment using immune system to fight cancer. Types: checkpoint inhibitors (PD-1, PD-L1, CTLA-4), CAR T-cell therapy, cytokines, cancer vaccines. Examples: pembrolizumab, nivolumab, ipilimumab, axicabtagene ciloleucel.

\medterm{Targeted Therapy}-- Drugs targeting specific molecular abnormalities in cancer. Examples: imatinib (BCR-ABL), trastuzumab (HER2), rituximab (CD20), tyrosine kinase inhibitors. Requires molecular testing.

\medterm{Hormone Therapy}-- Treatment blocking hormone effects in hormone-sensitive cancers. Examples: tamoxifen (breast cancer), leuprolide (prostate cancer), aromatase inhibitors.

\medterm{Tumor Marker}-- Substance produced by cancer or body in response to cancer. Used for diagnosis, monitoring, prognosis. Examples: PSA (prostate), CA-125 (ovarian), CEA (colon), AFP (liver), CA 19-9 (pancreas).

\medterm{Paraneoplastic Syndrome}-- Symptoms caused by cancer but not by tumor mass or metastasis. Mechanisms: hormone production, immune cross-reaction. Examples: SIADH, Cushing syndrome, Lambert-Eaton myasthenia, hypercalcemia.

\medterm{Cachexia}-- Wasting syndrome in cancer patients. Features: weight loss, muscle atrophy, anorexia, fatigue. Mediated by inflammatory cytokines. Difficult to treat, affects prognosis.

\medterm{Oncogene}-- Gene that can cause cancer when mutated or overexpressed. Normal version: proto-oncogene. Examples: RAS, MYC, HER2, BCR-ABL. Drive cell proliferation.

\medterm{Tumor Suppressor Gene}-- Gene that prevents uncontrolled cell growth. Examples: TP53, RB1, BRCA1, BRCA2, APC. Loss of function contributes to cancer.

\medterm{Apoptosis}-- Programmed cell death. Dysregulated in cancer (avoiding death). Pro-apoptotic genes: BAX, P53. Anti-apoptotic: BCL-2.

\medterm{Angiogenesis}-- Formation of new blood vessels. Tumors induce angiogenesis for growth. Targeted by anti-angiogenic drugs (bevacizumab).

\medterm{Metastatic Cascade}-- Sequence of events in metastasis: local invasion, intravasation, circulation, extravasation, colonization. Many cells die during process.

\medterm{Screening}-- Detection of cancer in asymptomatic individuals. Methods: mammography (breast), colonoscopy (colon), Pap smear (cervix), PSA (prostate), low-dose CT (lung). Balance benefits vs. harms.

\medterm{Prognosis}-- Predicted outcome of cancer. Factors: stage, grade, molecular features, patient status. Survivability statistics.

\medterm{Remission}-- Decrease or disappearance of cancer signs. Complete remission: no detectable disease. Partial remission: significant reduction but not complete.

\medterm{Recurrence}-- Return of cancer after treatment. Can be local, regional, or distant.

\medterm{Refractory}-- Cancer not responding to treatment. May require alternative approaches.

\medterm{Relapse}-- See recurrence.

\medterm{Primary Tumor}-- Original cancer site. Secondary: metastatic sites.

\medterm{Secondary Tumor}-- See metastasis.

\medterm{Primary Prevention}-- Preventing cancer development. Strategies: smoking cessation, vaccination (HPV, hepatitis B), sun protection, healthy diet, exercise.

\medterm{Secondary Prevention}-- Early detection through screening.

\medterm{Tertiary Prevention}-- Reducing complications from established cancer.

\medterm{Survival Rate}-- Percentage of patients surviving for specified time after diagnosis. 5-year survival commonly used.

\medterm{Palliative Care}-- Care focusing on symptom relief and quality of life. Can be provided with curative treatment.

\medterm{Hospice}-- End-of-life care for patients with <6 month prognosis. Comfort-focused.

\medterm{Clinical Trial}-- Research study testing new cancer treatments. Phases: I (safety), II (efficacy), III (comparison), IV (post-marketing).

\medterm{Informed Consent}-- Patient understanding and agreement to treatment or research participation.

\medterm{Comorbidity}-- See general.

\medterm{Adjuvant Therapy}-- Additional treatment after primary treatment to reduce recurrence risk. Examples: chemotherapy after surgery.

\medterm{Neoadjuvant Therapy}-- Treatment before primary treatment to shrink tumor. Examples: chemotherapy before surgery.

\medterm{Surgery}-- Surgical removal of tumor. Curative when complete resection possible. Palliative for symptom relief.

\medterm{Lymph Node Dissection}-- Removal of lymph nodes for staging or treatment. Sentinel node biopsy less invasive alternative.

\medterm{Sentinel Lymph Node}-- First lymph node draining tumor site. Biopsy assesses nodal spread.

\medterm{Imaging}-- See radiology.

\medterm{PET Scan}-- See nuclear medicine.

\medterm{MRI}-- See radiology.

\medterm{CT Scan}-- See radiology.

\medterm{Ultrasound}-- See radiology.

\medterm{Pathology}-- Study of disease. Cancer pathology: diagnosis, grading, staging, molecular characterization.

\medterm{Histology}-- Microscopic examination of tissue.

\medterm{Cytology}-- Examination of individual cells. Examples: Pap smear, fine needle aspiration.

\medterm{Immunohistochemistry}-- Use of antibodies to detect antigens in tissue. Used for cancer diagnosis and subtyping.

\medterm{Flow Cytometry}-- Analysis of cells by flow. Used for leukemia/lymphoma classification.

\medterm{Cytogenetics}-- Study of chromosomes. Detects abnormalities in cancer.

\medterm{Molecular Testing}-- Analysis of DNA/RNA for mutations, fusions, expression. Guides targeted therapy.

\medterm{Liquid Biopsy}-- Detection of circulating tumor DNA or cells in blood. Minimally invasive monitoring.

\medterm{Biomarker}-- Measurable indicator of disease. Cancer biomarkers: diagnostic, prognostic, predictive.

\medterm{Tumor Microenvironment}-- Surrounding tissue influencing tumor growth. Includes immune cells, blood vessels, fibroblasts, signaling molecules.

\medterm{Immune Evasion}-- Cancer mechanisms to avoid immune detection. Examples: PD-L1 expression, antigen loss.

\medterm{Checkpoints}-- Immune regulatory molecules. Cancer exploits checkpoints. Inhibitors block checkpoints.

\medterm{Checkpoint Inhibitor}-- Drug blocking immune checkpoints. Examples: pembrolizumab (PD-1), ipilimumab (CTLA-4).

\medterm{CAR T-Cell}-- Chimeric antigen receptor T-cells. Engineered immune cells targeting cancer.

\medterm{Monoclonal Antibody}-- Laboratory-produced antibody targeting specific antigen. Examples: rituximab, trastuzumab, bevacizumab.

\medterm{Vaccine}-- Preventive or therapeutic vaccination. Preventive: HPV vaccine, hepatitis B vaccine. Therapeutic: sipuleucel-T (prostate cancer).

\medterm{Hormone-Receptor Positive}-- Cancer expressing hormone receptors (ER+, PR+, AR+). Treated with hormone therapy.

\medterm{HER2-Positive}-- Cancer overexpressing HER2 protein. Treated with trastuzumab.

\medterm{Triple-Negative Breast Cancer}-- Breast cancer lacking ER, PR, HER2. Aggressive, limited treatment options.

\medterm{Small Cell Lung Cancer}-- Aggressive lung cancer, strongly associated with smoking. Early metastasis, responds to chemotherapy.

\medterm{Non-Small Cell Lung Cancer}-- Most common lung cancer type. Includes adenocarcinoma, squamous cell, large cell.

\medterm{Breast Cancer}-- Most common cancer in women. Risk factors: age, genetics, hormones, lifestyle. Screening: mammography. Treatment: surgery, radiation, chemotherapy, hormone therapy, targeted therapy.

\medterm{Lung Cancer}-- Leading cancer death worldwide. Types: small cell, non-small cell. Risk factors: smoking, radon, asbestos.

\medterm{Colorectal Cancer}-- Third most common cancer. Screening: colonoscopy. Prevention: polyp removal. Treatment: surgery, chemotherapy, radiation.

\medterm{Prostate Cancer}-- Common cancer in men. Screening: PSA, DRE. Most indolent, active surveillance often appropriate.

\medterm{Skin Cancer}-- Most common cancer. Types: basal cell, squamous cell, melanoma. Prevention: sun protection. Early detection critical.

\medterm{Melanoma}-- Most dangerous skin cancer. Risk factors: UV exposure, fair skin, moles. Early detection improves outcomes.

\medterm{Liver Cancer}-- Hepatocellular carcinoma, often in cirrhosis. Risk factors: hepatitis, alcohol, aflatoxin.

\medterm{Pancreatic Cancer}-- Aggressive cancer, often diagnosed late. Risk factors: smoking, chronic pancreatitis, genetics.

\medterm{Kidney Cancer}-- Renal cell carcinoma. Risk factors: smoking, obesity, hypertension.

\medterm{Bladder Cancer}-- Urothelial carcinoma. Risk factors: smoking, occupational exposures.

\medterm{Thyroid Cancer}-- Usually indolent. Types: papillary, follicular, medullary, anaplastic.

\medterm{Leukemia}-- See leukemia.

\medterm{Lymphoma}-- See lymphoma.

\medterm{Myeloma}-- Multiple myeloma: plasma cell malignancy. Features: bone pain, anemia, renal failure, hypercalcemia.

\medterm{Sarcoma}-- See sarcoma.

\medterm{Pediatric Cancer}-- Cancers in children. Common: leukemia, brain tumors, lymphoma, neuroblastoma, Wilms tumor. Often curable.

\medterm{Survivorship}-- Care for cancer survivors. Addresses late effects, surveillance, quality of life.
